\begin{table}[htbp]
\centering
\caption{Result 3: Does the Premium Scale with the Nonlegal-Side Population?}
\label{tab:result3}
\small
\begin{tabular}{p{7.4cm}cccc}
\toprule
Specification & Coef. & SE (HC1) & SE (cl.) & $p$ (cl.) \\
\midrule
log($P_i$: captive nonlegal population, 70mi) + full controls & $+0.030$ & $0.007$ & $0.014$ & $0.036$ \\
\quad excl.\ Sedgwick County, CO & $+0.024$ & $0.007$ & $0.014$ & $0.078$ \\
\bottomrule
\end{tabular}
\begin{minipage}{0.95\textwidth}
\vspace{0.3em}
\footnotesize
\textit{Notes:} Outcome is log(revenue per active dispensary), $n=787$. $P_i$ sums the population of
\emph{all} nonlegal-state counties within 70 miles of county $i$, not only the nearest one, so it
distinguishes a county next to a small nonlegal county from a county next to a large nonlegal metropolitan
area, something the distance measures in Result 1 cannot do. Only 10 of the 80 Colorado and Washington
counties in this panel have any captive nonlegal population within 70 miles at all, so this result reweights
the same identifying set of border counties from Result 1 by market size rather than adding new geographic
variation; combined with Result 1's continuous-distance finding, it indicates the premium responds
continuously to both how close a county is to a nonlegal border and how large the population on the other
side is, not only to a binary border classification.
\end{minipage}
\end{table}
